\chapter{Implementation}
\label{ch:implementation}

This chapter describes the implementation of Ai4k8s, focusing on the engineering
challenges encountered during development and the solutions adopted. Section~\ref{sec:infra}
covers the infrastructure and deployment configuration. Section~\ref{sec:mubench-impl}
describes the muBench workload setup. Section~\ref{sec:gaps} documents the six
engineering gaps identified during evaluation and their resolutions.
Section~\ref{sec:topsis-weights} presents the TOPSIS weight calibration study.

% -----------------------------------------------------------------------
\section{Infrastructure and Deployment}
\label{sec:infra}

All experiments were conducted on a CrownLabs virtual machine instance with the
specifications in Table~\ref{tab:vm-spec}.

\begin{table}[h]
  \centering
  \caption{CrownLabs VM specifications.}
  \label{tab:vm-spec}
  \begin{tabular}{ll}
    \toprule
    Parameter & Value \\
    \midrule
    Instance ID    & instance-866fd \\
    vCPUs          & 4 \\
    RAM            & 14~GiB \\
    Disk           & 42~GB \\
    OS             & Ubuntu 22.04 \\
    Kubernetes     & k3s v1.34.5+k3s1 (single-node) \\
    Ollama         & v0.17.7 \\
    LLM model      & qwen3:0.6b (Q4, 522~MB) \\
    \bottomrule
  \end{tabular}
\end{table}

\subsection{Service Architecture}
Three systemd services manage the runtime components:
\begin{itemize}
  \item \textbf{ai4k8s-web.service} — Flask+SocketIO web application (port 5003).
  \item \textbf{mcp-http.service} — FastAPI MCP server (port 5002, localhost only).
  \item \textbf{ollama.service} — Ollama LLM server (port 11434, localhost only).
\end{itemize}

\subsection{Disk Pressure Mitigation}
Ollama installs GPU runtime libraries (CUDA v12/v13, Vulkan, MLX) by default
despite the VM having no GPU. These libraries occupy 4.3~GB of disk. On a VM
with a 15~GB disk (as with the previous 7.2~GiB RAM instance), this consumed
enough space to trigger a Kubernetes \textit{disk-pressure} taint, blocking all
pod scheduling. The libraries were removed immediately after Ollama installation:

\begin{verbatim}
sudo rm -rf /usr/local/lib/ollama/cuda_v12 \
            /usr/local/lib/ollama/cuda_v13 \
            /usr/local/lib/ollama/mlx_cuda_v13 \
            /usr/local/lib/ollama/vulkan
\end{verbatim}

This is now part of the instance setup checklist and reduces disk usage from
$\sim$6.8~GB to $\sim$2.5~GB post-install.

% -----------------------------------------------------------------------
\section{muBench Microservice Workload}
\label{sec:mubench-impl}

The muBench framework~\cite{detti2021mubench} provides a configurable
microservice benchmarking workload. Each service is a Python HTTP server that
executes a user-defined compute function when it receives a request, optionally
forwarding downstream calls to other services in a configurable call graph.

\subsection{Service Chain Configuration}
We deploy a three-service chain: \textit{ingest} $\to$ \textit{process}
$\to$ \textit{analyze}. Table~\ref{tab:mubench-services} describes each service.

\begin{table}[h]
  \centering
  \caption{muBench service chain configuration.}
  \label{tab:mubench-services}
  \begin{tabular}{llll}
    \toprule
    Service & Compute function & CPU limit & RAM limit \\
    \midrule
    ingest  & image\_processing (convolution + histogram EQ) & 500m & 256~MiB \\
    process & matrix\_compute (4-layer YOLO-style forward pass) & 500m & 256~MiB \\
    analyze & video\_frames (optical flow + FFT energy) & 300m & 256~MiB \\
    \bottomrule
  \end{tabular}
\end{table}

Each service is annotated for Ai4k8s:
\begin{verbatim}
kubectl annotate deployment ingest \
  ai4k8s.io/predictive-autoscaling-enabled=true \
  autosage.ai4k8s/state-management=stateless
\end{verbatim}

\subsection{seq\_len Parameter}
\label{sec:seqlen}

The \texttt{seq\_len} field in \texttt{workmodel.json} controls the fan-out:
each upstream request generates \texttt{seq\_len} downstream calls.

With \texttt{seq\_len=100}: one user request generates $100 \times 100 = 10{,}000$
\textit{analyze} calls, saturating all three services simultaneously. This makes
it impossible to demonstrate per-service differentiated scaling because all HPAs
hit \texttt{maxReplicas} within 90~seconds of each other.

Reducing to \texttt{seq\_len=5} produces $5 \times 5 = 25$ \textit{analyze}
calls per user request. Under 32 concurrent connections, \textit{ingest} saturates
first at $\sim$60\% CPU while \textit{process} and \textit{analyze} remain below
15\%, demonstrating the intended bottleneck-first scaling story.

\subsection{Deployment Challenges}
Several underdocumented muBench requirements were discovered during deployment:
\begin{itemize}
  \item The container entrypoint requires five environment variables
  (\texttt{APP}, \texttt{ZONE}, \texttt{K8S\_APP}, \texttt{PN}, \texttt{TN})
  that are not in the official documentation; found by reading the controller source.
  \item The default liveness probe path (\texttt{/api/v1}) executes the full
  compute chain, causing timeouts; replaced with \texttt{/metrics}.
  \item The workmodel call graph required explicit port numbers in service URLs;
  absent ports caused silent connection timeouts.
\end{itemize}

% -----------------------------------------------------------------------
\section{Engineering Gaps and Resolutions}
\label{sec:gaps}

Six engineering gaps were identified between the initial design and a fully
functional evaluation. Table~\ref{tab:gaps} summarizes the gaps and fixes.

\begin{table}[h]
  \centering
  \caption{Engineering gaps identified and resolved during implementation.}
  \label{tab:gaps}
  \small
  \begin{tabular}{clll}
    \toprule
    ID & Description & Root cause & Fix \\
    \midrule
    1A & Ollama blocks cascade & Native call hangs 240~s & 90~s soft timeout thread \\
    1B & Groq not triggered    & TimeoutError swallowed & Propagate exception correctly \\
    1C & Sequential loop (3$\times$ slower) & ThreadPoolExecutor missing & Parallel per-service decisions \\
    2  & force\_rising not firing & Missing eval flag & \texttt{force\_rising=True} in eval script \\
    3  & VPA poll window too short & 120~s $<$ recommender convergence & \texttt{VPA\_POLL\_WINDOW=300~s} \\
    4  & Cached LLM responses & Session cache not cleared & Inter-cycle cache flush \\
    5  & TOPSIS gap=0 always   & forecast weight=0.15 too low & Raised to 0.25 \\
    6  & readyReplicas crash   & Null during rollout & Fallback to \texttt{spec.replicas} \\
    \bottomrule
  \end{tabular}
\end{table}

\subsection{Gap 1A — Ollama Soft Timeout}
\label{sec:gap1a}

The original implementation called Ollama synchronously in the main autoscaling
loop. Under CPU-saturating load, GGUF inference can take over 240~seconds,
blocking the entire background thread and preventing any autoscaling action from
completing. A thread-based soft timeout was introduced:

\begin{algorithm}
  \caption{Ollama soft timeout wrapper}
  \label{alg:soft-timeout}
  \begin{algorithmic}[1]
    \State $result \gets \{\}$
    \State \textbf{def} $\_ollama\_call()$: $result \gets$ \textbf{call} Ollama API
    \State $t \gets$ Thread(target=$\_ollama\_call$, daemon=True)
    \State $t$.start(); $t$.join(timeout=$T_{soft}$)
    \If{$result$ is set}
      \State \Return $result$
    \Else
      \State \textbf{raise} TimeoutError($T_{soft}$)
    \EndIf
  \end{algorithmic}
\end{algorithm}

The daemon thread continues running after the timeout (it cannot be killed cleanly),
but the main loop proceeds to Groq fallback. $T_{soft} = 90$~s was selected as a
balance between giving the local model time to respond at high load and limiting
the delay before cloud fallback.

\subsection{Gap 1C — Parallel Background Loop}
\label{sec:gap1c}

The initial background loop processed the three muBench services sequentially.
With Qwen Q8\_0 at 220~s average inference, processing all three services took
$\sim$660~s—longer than the 5-minute autoscaling interval. A
\texttt{ThreadPoolExecutor} with \texttt{max\_workers=3} was introduced to
dispatch all three service decisions in parallel:

\begin{verbatim}
with ThreadPoolExecutor(max_workers=3) as executor:
    futures = {
        executor.submit(predict_and_scale, svc): svc
        for svc in monitored_services
    }
    for future in as_completed(futures):
        results[futures[future]] = future.result()
\end{verbatim}

With qwen3:0.6b and Groq fallback, all three per-cycle decisions complete within
5--6~seconds of wall time.

\subsection{Gap 4 — LLM Response Cache}
\label{sec:gap4}

Ai4k8s caches LLM responses keyed on deployment name and CPU value to avoid
redundant inference within a single web session. The background evaluation loop
reuses this cache across cycles if metrics are unchanged, producing identical
decisions without querying the LLM. The fix clears the cache at the start of
each evaluation cycle:

\begin{verbatim}
self.recommendation_cache.clear()
\end{verbatim}

This ensures every cycle triggers a fresh LLM inference and captures authentic
decision latency data.

\subsection{Gap 5 — TOPSIS Weight Calibration}
\label{sec:weight-study}

With the initial weights (forecast\_alignment=0.15), the TOPSIS score gap between
\texttt{scale\_up} and \texttt{maintain} was 0.2585 for \textit{ingest} (49.1\%
CPU, rising forecast). This gap indicates MCDA overrode the LLM recommendation
but the margin was marginal. Raising forecast\_alignment to 0.25 (and reducing
cost from 0.20 to 0.15) collapsed the gap to 0.0000 in both directions—the
criteria now balance such that the LLM forecast and the stability/cost
constraints counterbalance cleanly, producing full agreement (gap=0.0000)
in the stable-load evaluation.

% -----------------------------------------------------------------------
\section{Qwen Think-Mode Bypass}
\label{sec:think-bypass}

Qwen3 models use chain-of-thought by default. When called through the
OpenAI-compatible endpoint (\texttt{/v1/chat/completions}), the thinking tokens
are consumed internally and the returned \texttt{content} field is empty.
Ai4k8s detects this by checking whether the configured base URL contains
\texttt{:11434} (Ollama's default port) and routes those requests through the
native \texttt{/api/chat} endpoint with \texttt{"think": false}:

\begin{verbatim}
if ":11434" in self.base_url:
    # Use native Ollama endpoint to disable chain-of-thought
    response = requests.post(
        f"{self.base_url.replace('/v1', '')}/api/chat",
        json={"model": self.model, "think": False, ...}
    )
\end{verbatim}

This workaround is transparent to the rest of the pipeline: the native and
OpenAI-compatible paths return identically structured response objects.
